\documentclass[11pt]{article}
\usepackage[T1]{fontenc}
\usepackage[utf8]{inputenc}
\usepackage{lmodern}
\usepackage[margin=1in]{geometry}
\usepackage{amsmath,amssymb,amsthm,mathtools}
\usepackage{booktabs,array,longtable}
\usepackage{microtype}
\usepackage{enumitem}
\usepackage{listings}
\usepackage[hidelinks]{hyperref}
\usepackage{fancyhdr}
\setlength{\headheight}{14pt}
\pagestyle{fancy}
\fancyhf{}\fancyhead[L]{\small Exact acceleration of a scanning unknot recognizer}
\fancyhead[R]{\small September 2026}\fancyfoot[C]{\thepage}
\newtheorem{theorem}{Theorem}[section]
\newtheorem{proposition}[theorem]{Proposition}
\newtheorem{lemma}[theorem]{Lemma}
\newtheorem{corollary}[theorem]{Corollary}
\theoremstyle{definition}\newtheorem{definition}[theorem]{Definition}
\theoremstyle{remark}\newtheorem{remark}[theorem]{Remark}
\newcommand{\F}{\mathbb F_2}
\newcommand{\Kh}{\operatorname{Kh}}
\newcommand{\rKh}{\widetilde{\operatorname{Kh}}}
\newcommand{\rk}{\operatorname{rank}}
\newcommand{\End}{\operatorname{End}}
\newcommand{\im}{\operatorname{im}}
\newcommand{\poly}{\operatorname{poly}}
\newcommand{\unk}{\textsf{UNKNOT}}
\newcommand{\knot}{\textsf{KNOTTED}}
\newcommand{\unknown}{\textsf{UNKNOWN}}
\lstset{basicstyle=\small\ttfamily,columns=fullflexible,keepspaces=true,
  breaklines=true,showstringspaces=false,frame=single,framerule=.3pt,
  xleftmargin=3pt,xrightmargin=3pt}
\title{Exact Acceleration of a Scanning Unknot Recognizer\\[4pt]
\large Sparse cancellation, finite-ring obstructions,\\and connected-sum compression}
\author{Computational research and implementation report\\
\small Companion to the accelerated \texttt{fastunknot} source distribution}
\date{18 September 2026}
\begin{document}
\maketitle
\begin{abstract}
We improve the Python unknot recognizer in the supplied \texttt{Knots.zip},
rather than replacing its complete algorithm with an incomplete invariant test.
The implementation combines exact one-sided modular Alexander and Jones
obstructions, a faster descending-diagram test, an asymptotically faster version
of the original greedy scan ordering, sparse unit cancellation in the dotted
cobordism category, and visible connected-sum factorization. Every successful
verdict retains a complete mathematical justification; resource exhaustion is
reported separately. In process-isolated tests, complete recognition of the
Conway and Kinoshita--Terasaka examples improves by factors of about 33 and 24.
An unfactored 36-crossing backend workload that did not finish in a separate
120-second baseline run completes in 2.81 seconds, with reduced rank 2,949.
For connected sums of fixed bounded-size diagrams, factorwise rank computation
is polynomial, while the original explicit minimal-complex representation can
require exponentially many objects even at constant boundary width. A
132-crossing example with reduced rank $33^{12}$ takes 0.18 seconds and a peak of
246 objects in any factor scan. We prove the transformations, describe the
validation and measured regressions, and audit the relationship to Lackenby's
hierarchy approach. \emph{No general quasi-polynomial bound is established for
the delivered implementation.} Its general worst-case bound remains
$2^{O(n)}$; the asymptotic gains are in preprocessing and specified decomposable
families, in addition to the measured practical gains.
\end{abstract}
\tableofcontents

\section{Scope, input, and the baseline}
The target is the \texttt{fast/} implementation in the user's archive. Its
unmodified Python package, tests, and examples are retained under
\texttt{baseline/}; the replacement package is \texttt{fastunknot/}. The
result is executable with the Python standard library alone. The report is
about this implementation and these experiments, not a claim to have developed
a new mathematical unknot-detection theorem.

An input is a classical, one-component knot diagram with $n$ crossings. A planar
diagram (PD) row is $[a,b,c,d]$ in counterclockwise order, with the under-strand
joining $a$ to $c$ and the over-strand joining $b$ to $d$. Edge labels occur twice
and are normalized to $0,\ldots,2n-1$. The empty PD denotes \emph{one} circle,
not the empty link. Constructors also accept braid closures and rectangular
(grid) diagrams. They check the strand-component count and the Euler
characteristic of the specified rotation system. An abstractly planar graph
with a nonspherical rotation system is not accepted as a classical diagram.
Low-level scan routines retain the baseline precondition that their PD is valid;
applications should use the validated \texttt{Diagram} constructors.

The original recognition pipeline performs crossing-decreasing Reidemeister I
and II moves, tests for a descending traversal, computes an exact Alexander
polynomial, and finally computes Khovanov homology over $\F$ by scanning. The
last step is based on local complexes, delooping, and Gaussian cancellation in
the dotted-cobordism framework of Bar--Natan~\cite{BN}. The implementation forgets
quantum degree and retains unnormalized cube degree. For a knot, the total
unreduced $\F$ rank is twice the reduced rank. Reduced rank one detects the
unknot; Section~\ref{sec:complete} explains the coefficient-field issue.

The supplied baseline is already a genuine complete recognizer when resource
limits are absent. The improvement must therefore preserve exactness, not merely
recognize a larger collection of easy examples. Our changes fall into three
categories: polynomial preprocessing improvements, exact rejection filters, and
changes to the computation and representation of the complete homological
fallback. The raw-backend and complete-recognition timings are deliberately
reported separately.

\section{Attempting the quasi-polynomial route}\label{sec:quasi}
The attached talk announces a quasi-polynomial algorithm. Its final flowchart
contains hierarchical multi-surfaces, cutting, Cheeger-region simplification,
Haken's lemma, weak reduction, and multi-surface simplification. The preceding
slides use generalized Heegaard splittings and explain that absence of a
Cheeger region yields logarithmic hierarchy length~\cite{Ltalk}. These are
three-dimensional topological operations, not alternate names for Khovanov
matrix pivots.

The attached July 2026 paper is a related hierarchy algorithm. Section 9 gives
an iteration bound $L(g+1)^L$ and explicitly distinguishes it from a
running-time estimate. Section 10 then refines the algorithm using interior
parallelity bundles and gives a linear bound on $L$ in the initial
$q$-complexity~\cite{L2026}. Thus it would be inaccurate either to call the
Section 9 formula a proved bit-complexity bound, or to say that the paper never
controls hierarchy length.

\subsection{The quantitative bridge that an implementation would need}
Let $B(n)$ bound the number of bits in a maintained topological state and let
$P(B)$ bound the cost of an elementary operation. Suppose a hierarchy-based
implementation established
\begin{equation}\label{eq:hierbound}
 T(n)\leq n^a L(g+1)^L P(B(n)),\qquad
 L\leq c\log(n+2),\qquad g\leq(n+2)^d,
\end{equation}
with $B(n)\leq n^{O(1)}$ and polynomial $P$. Taking logarithms then gives
\[
 \log T(n)=O(\log n)+O(\log n)\,O(\log n)=O((\log n)^2),
\]
and hence $T(n)=n^{O(\log n)}$. This elementary calculation identifies what
would actually have to be measured and proved. Replacing the logarithmic bound
on $L$ by a linear bound in $n$ gives only
$L(g+1)^L=\exp(O(n\log n))$ under the same polynomial bound on $g$.
Furthermore, a small \emph{number} of transitions does not prevent a transition
from expanding an exponentially long normal surface.

A viable implementation contract would therefore need a compressed state
containing the manifold's handle incidences, boundary pattern, surface
coordinates, hierarchy references, and Heegaard data, together with actual
algorithms for the transformations above. Each transformation must preserve the
represented topology, maintain explicit bit-size bounds, and either decrease
a proved potential or make progress in a bounded construction loop. Merely
creating stubs named after the flowchart would not meet that contract.

The supplied scanning objects do not encode this information. In particular,
their boundary matching is a two-dimensional smoothing interface; it is not a
normal surface in a knot exterior. I did not obtain implementations and proofs
of these three-dimensional contracts in this work. The delivered executable
therefore does not call a putative quasi-polynomial subroutine or advertise
such a guarantee. Instead, the following sections give completed improvements
with explicit correctness arguments and executed tests. This is a conclusion
about the present implementation, not a claim that the hierarchy program is
impossible or that the announced theorem is false.

\section{The exact portfolio and its decision invariant}\label{sec:portfolio}
The accelerated recognizer follows this sequence:
\begin{center}
\begin{minipage}{0.94\linewidth}
\begin{enumerate}[itemsep=2pt,topsep=3pt]
\item Apply the original validated decreasing R1/R2 moves; an empty diagram is $\unk$.
\item Find a descending traversal by cyclic intervals; success gives $\unk$.
\item Split visible connected sums. A nontrivial factor gives $\knot$;
      all trivial factors give $\unk$.
\item Apply modular Alexander necessary conditions; a mismatch gives $\knot$.
\item Evaluate a normalized bracket modulo a fixed prime; a mismatch gives $\knot$.
\item Compute the original exact Alexander polynomial; a nonunit gives $\knot$.
\item Compute reduced $\F$ Khovanov rank with accelerated scanning; rank one
      gives $\unk$ and any other rank gives $\knot$.
\end{enumerate}
\end{minipage}
\end{center}
An inconclusive filter proceeds to the next stage. Exceeding the optional Jones
state cap also proceeds to the next stage; exceeding the overall user time or
Khovanov-object budget produces $\unknown$. No conclusion is inferred from an
exhausted search, a modular match, or the failure to find a decreasing move.

The default Jones state cap is 20,000. Modular Alexander uses the fixed prime
$p=1{,}000{,}000{,}007$, verified in the tests by trial division through its
square root. Neither numerical roundoff nor a probabilistic identity test enters
a verdict. Although a modular evaluation can miss a nontrivial polynomial,
that can only cause more work, not a false $\unk$ answer.

\section{Faster preprocessing without changing its meaning}
\subsection{Linear descending-start detection}
Choose one orientation and write the $2n$ incoming darts encountered around the
knot as a cyclic sequence. For crossing $c$, let $o_c$ and $u_c$ be the indices
of its over-visit and under-visit. Starting the traversal at index $s$ encounters
$c$ under first exactly when
\begin{equation}\label{eq:badinterval}
 s\in(o_c,u_c]\quad\text{in cyclic order}.
\end{equation}
A descending start therefore lies outside the union of these forbidden
intervals.

\begin{proposition}
A difference-array computation finds a descending start, or certifies that none
exists, in $O(n)$ combinatorial operations and $O(n)$ storage.
\end{proposition}
\begin{proof}
Split a wrapping cyclic interval into at most two ordinary intervals. For an
interval $[l,r]$, add one at $l$ and subtract one at $r+1$. Prefix sums count
forbidden intervals covering each candidate start. A zero count is precisely
a start that encounters every crossing over first, by~\eqref{eq:badinterval}.
Perform the same calculation in the opposite orientation. If $d$ is an
incoming dart, the opposite incoming dart is the port two positions away at
the same crossing; reversing the traversal and applying this port involution
gives the other orientation. Thus the two linear scans cover the same $4n$
possible starting darts examined separately by the baseline.
\end{proof}
The original implementation has an $O(n^2)$ worst-case upper bound because it
restarts a traversal for every candidate. Some inputs fail almost immediately
at each candidate, so their actual baseline cost can already be linear. The
501-crossing microbenchmark exhibits essentially no practical gain; this does
not contradict the improved worst-case preprocessing bound.

\subsection{The same greedy order in near-linear time}
At any partial scan, let $s_i$ be the number of edges joining an unprocessed
crossing $i$ to the processed set, and let $\ell_i$ be its number of self-loops.
The original greedy score, after removing quantities common to all candidates,
orders crossings lexicographically by
\[
 (s_i,\ell_i,-i).
\]
Indeed, if the current boundary has size $b$, processing $i$ changes it to
$b+4-2s_i-2\ell_i$. The tie-breaking rule is unchanged.

Store $(-s_i,-\ell_i,i)$ in a min-heap. After processing a crossing, only its
at most four incident non-loop neighbours can acquire additional shared edges.
Insert their new scores and discard stale heap entries lazily. The total number
of insertions is $O(n)$ and heap operations take $O(\log n)$ each. This yields
the exact same order for a specified starting crossing as the original
$O(n^2)$ scan over all remaining candidates.

The default tries a bounded number of starts. The inherited spacing rule can
try somewhat more than 12 starts, but fewer than 24; thus the default order
selection is $O(n\log n)$, versus $O(n^2)$ previously. With every crossing used
as a start, the bounds become $O(n^2\log n)$ and $O(n^3)$. These bounds count
operations on normalized labels; bit costs and Python allocation constants
are additional implementation details. Five hundred randomized comparisons
check equality of the old and new orders at up to three starts each.

\section{Exact finite-ring obstructions}\label{sec:filters}
\subsection{Evaluating an Alexander minor directly}
Use the same Wirtinger/Fox rows as the original package. Deleting one row and
one column yields a polynomial determinant $D(t)$ representing the Alexander
polynomial up to a unit $\pm t^k$. Rather than construct polynomial coefficient
lists, evaluate each sparse row at $t=a\in\mathbb F_p$ first and perform ordinary
Gaussian elimination over $\mathbb F_p$. This costs $O(n^3)$ field operations;
$p$ is fixed, so all filter coefficients have bounded bit length.

\begin{proposition}\label{prop:alexfilter}
For a knot with unit Alexander polynomial, the following necessary conditions
hold whenever $a\ne0$ in $\mathbb F_p$:
\[
 D(-1)\in\{1,-1\},\qquad D(a)D(a^{-1})=1.
\]
Consequently, violating either condition proves nontriviality.
\end{proposition}
\begin{proof}
For a unit polynomial, $D(t)=\varepsilon t^k$ with
$\varepsilon\in\{1,-1\}$. Substitution gives
$D(-1)=\varepsilon(-1)^k$ and
$D(a)D(a^{-1})=\varepsilon^2 a^ka^{-k}=1$. Evaluation is a ring
homomorphism, so computing the determinant after evaluation gives the same
residue as evaluating the determinant. The unknot has unit Alexander
polynomial, establishing the contrapositive test.
\end{proof}
The implementation tries $t=-1$ first. If inconclusive, it tries $a=2$ and
$a^{-1}$. The reciprocal product avoids guessing the unknown sign and power
of $t$. It also catches some nontrivial Alexander polynomials with determinant
one. It does not catch the Conway or Kinoshita--Terasaka examples, whose
Alexander polynomial is identically one. The exact polynomial stage remains
available after the modular and Jones stages.

\subsection{A writhe convention that must not be guessed}
The original public \texttt{signs()} routine used the incoming over-port alone.
That value serves its fixed-port Fox-row convention, but is not the oriented
crossing sign on an arbitrary valid PD. Let $u_i\in\{0,2\}$ and
$o_i\in\{1,3\}$ be the incoming under- and over-ports. The sign compatible
with this package's bracket convention is
\begin{equation}\label{eq:sign}
 \sigma_i=
 \begin{cases}1&o_i=3,\\-1&o_i=1\end{cases}
 \begin{cases}1&u_i=0,\\-1&u_i=2.\end{cases}
 \qquad w(D)=\sum_i\sigma_i.
\end{equation}
A half-turn of a PD row changes both factors, leaving the product invariant.
Reversing the knot orientation likewise changes both incoming ports and leaves
$\sigma_i$ invariant. These two requirements immediately expose the one-port
problem. In the inherited braid-to-PD convention, the single generator
\texttt{[1]} has $w=-1$; callers must not impose an unrelated sign convention.

The old one-port values are retained as \texttt{relation\_signs()}, and the
Alexander code explicitly calls that routine. Its $a$-port is treated as the
incoming port in the chosen relation, whether or not it is globally incoming.
Simply replacing those relation signs with~\eqref{eq:sign} would introduce a
new error. The original recognition decisions did not use Jones normalization,
so this finding is not evidence that the baseline's Alexander verdicts were
incorrect. It is a public writhe bug and a prerequisite for a safe new filter.

\subsection{Loop-inclusive normalized bracket}
Fix a commutative ring and a unit $A$. Write
$\delta=-A^2-A^{-2}$. For a nonempty diagram define
\begin{equation}\label{eq:bracket}
 B_D(A)=\sum_{s\in\{0,1\}^n}A^{n-2|s|}\delta^{c(s)},
 \qquad F_D(A)=(-A^3)^{-w(D)}B_D(A),
\end{equation}
where $c(s)$ counts all closed smoothing circles. This is the
loop-inclusive form of the bracket state model~\cite{Kauffman}; in particular
$B_U=F_U=\delta$, rather than 1. The empty PD is handled as the one-circle
special case.

The smoothing connecting $(a,b),(c,d)$ has coefficient $A$, and the smoothing
connecting $(a,d),(b,c)$ has coefficient $A^{-1}$. The local R1 calculation
multiplies $B$ by $-A^3$ or $-A^{-3}$, exactly cancelled by the writhe factor
under~\eqref{eq:sign}. R2 follows from the local relation $e^2=\delta e$:
\[
 (A\,1+A^{-1}e)(A^{-1}\,1+Ae)
 =1+(A^2+A^{-2}+\delta)e=1.
\]
Expanding the two three-crossing products, using the adjacent smoothing
relations $e_i e_{i+1}e_i=e_i$, gives R3. Thus $F_D$ is a knot invariant with
unknot value $\delta$. These polynomial identities remain valid after
specialization in any ring where $A$ is invertible.

\begin{proposition}\label{prop:jonesfilter}
A modular value $F_D(A)\ne\delta$ proves that $D$ is knotted. Equality is
inconclusive, including when the chosen specialization makes $\delta=0$.
\end{proposition}
This test never assumes that the full Jones polynomial detects the unknot.
It also never assumes that one evaluation determines the polynomial. For
example, the trefoil at $A=1$ passes this filter; the tests check that this
collision is not interpreted as an unknot certificate.

\subsection{Matching-state dynamic programming}
A partial smoothing is represented by its matching of exposed edge endpoints.
For each matching, store a scalar coefficient. Attaching the next crossing has
two choices. Multiply by $A$ or $A^{-1}$ and by $\delta^c$ for the newly
closed circles, then add coefficients of equal resulting matchings. Induction
on the processed crossings proves that each stored coefficient is the sum of
exactly the partial states it represents. At the end only the empty matching
remains, with coefficient $B_D(A)$.

If the boundary has $w$ endpoints, an unconditional upper bound on the number
of possible matching states is
\[
 (w-1)!!=\frac{w!}{2^{w/2}(w/2)!}=2^{O(w\log w)}.
\]
A Catalan bound requires additional disk-boundary/noncrossing structure, which
we do not assume for an arbitrary scan prefix. The state count is also at most
the number of processed resolutions. With the fixed state cap, the optional
filter has bounded-state polynomial overhead; without a cap, no general
polynomial bound in $n$ is asserted. A cap hit discards the partial evaluation
and falls through to the complete recognizer.

\section{Accelerating exact sparse cancellation}\label{sec:cancellation}
\subsection{What the scanner stores}
A scanning object consists of a boundary matching and a cube degree. Closed
circles are delooped into choices of the two labels $1,x$ in
$\F[x]/(x^2)$. Between matchings $a,b$, the morphism representation is a set of
square-free dot monomials on the circles obtained by gluing $a$ to the reflected
$b$. Addition is symmetric difference. Composition is the exact two-dimensional
Frobenius-algebra calculation, not a numerical approximation.

The geometry routines of the supplied implementation are preserved. Most new
work removes repeated computation and chooses a different valid elimination
order. Local tensoring, delooping, and cancellation are the algebraic basis of
the fast scanning method~\cite{BN}. Forgetting quantum grading is sufficient for
total-rank detection, but these routines do not output integral torsion or a
bigraded polynomial.

\subsection{Every unit is self-inverse in characteristic two}
For a matching with $r$ arcs, its endomorphism ring in the implemented normal
form is
\[
 R_r=\F[x_1,\ldots,x_r]/(x_1^2,\ldots,x_r^2).
\]
A morphism is a unit precisely when its constant term is one.

\begin{lemma}\label{lem:inverse}
If $f=1+\nu\in R_r$ and $\nu$ has zero constant term, then $\nu^2=0$ and
$f^{-1}=f$.
\end{lemma}
\begin{proof}
Write $\nu$ as a sum of nonconstant square-free monomials. In characteristic
two the cross terms in its square cancel in pairs. The square of each remaining
monomial contains $x_i^2$ for at least one $i$ and is zero. Therefore
$(1+\nu)^2=1+2\nu+\nu^2=1$.
\end{proof}
The baseline computed a geometric series for the inverse. The replacement
returns a fresh copy of $f$. It is not merely an early truncation chosen from
observed data. The same proof explains a direct local-ring composition routine:
two dot monomials multiply to zero if they overlap and to their union otherwise;
equal output monomials cancel by symmetric difference. Identity and zero
compositions are also recognized before invoking geometry.

\subsection{Fill-aware Gaussian pivots}
Suppose an invertible differential entry $\phi:b\to c$ is selected. For
surviving entries $\delta:a\to c$ and $\gamma:b\to f$, cancellation changes
\begin{equation}\label{eq:schur}
 d_{af}\longleftarrow d_{af}+\gamma\phi^{-1}\delta.
\end{equation}
The sign is plus because the coefficient field is $\F$.

\begin{lemma}\label{lem:cancel}
Removing $b,c$ and updating by~\eqref{eq:schur} preserves the chain-homotopy
type and hence the homology rank.
\end{lemma}
\begin{proof}
In the relevant two degrees, isolate the invertible block $\phi$ in the
matrix of the differential. Invertible row and column operations clear its
row and column, leaving the Schur complement on the other objects. The cleared
block is the two-term contractible complex $b\xrightarrow{\phi}c$.
The equation $d^2=0$ makes the adjacent differential blocks compatible with
these changes of basis. Deleting the contractible summand gives exactly
\eqref{eq:schur}. The argument works in the additive category because the
only required inverse is that of $\phi$.
\end{proof}
The number of possible new entries from this cancellation is bounded by
\[
 \mu(b,c)=(|\operatorname{inc}(c)|-1)(|\operatorname{out}(b)|-1).
\]
The new implementation keeps unit entries in a lazy heap ordered by this
Markowitz-style fill estimate, with deterministic object-ID ties. It recomputes
a popped score and reinserts it when stale. Scores of other heap entries can
still be stale, so the implementation does \emph{not} promise a globally
minimum current score. What matters for correctness is that every performed
pivot is genuinely invertible. Every unit present initially is enqueued, and
every changed entry that becomes a unit is enqueued. Thus changing the queue
policy cannot stop elimination with an unexamined surviving unit.

There is no theorem here that this heuristic improves every input. Cancellation
order can change transient fill dramatically, but it can also introduce heap
overhead with little benefit. The experiments measure the complete package of
kernel changes, not a separately randomized causal attribution to each change.

\subsection{Caching without hidden algebraic mutation}
Composition is bilinear. We cache the composition of individual input
monomials and sum the immutable cached answers. Crossing-tensor entries are
also cached with immutable morphism values; callers receive fresh mutable
containers. The geometry caches use 16,384 entries each, the monomial
composition cache 32,768, and the crossing-entry cache 8,192. Cache eviction
only forces recomputation and has no mathematical effect.

Entry-count bounds are not byte bounds: one morphism can itself be large.
The internal geometry cache objects must be treated as read-only. The tests
exercise mutation of returned morphism results to check that a caller cannot
corrupt a cached composition through the public arithmetic wrapper. The
\texttt{clear\_caches()} function supports cold experiments and long-lived
applications. These caches are performance aids, not persistent mathematical
assumptions.

\section{Connected sums and avoiding explicit homology multiplicities}
\label{sec:factor}
\subsection{Finding visible splitting curves}
In the planar dual of the knot projection, two parallel edges joining distinct
faces form a simple two-edge cycle. Its geometric realization crosses exactly
two edges of the original diagram. The Jordan curve separates the sphere into
two disks, and closing the strand segment on each side produces diagrams
$D_1,D_2$ with $D=D_1\#D_2$.

The implementation computes dart-to-face incidences and groups primal edges by
the unordered pair of their incident faces. For a candidate pair it deletes the
two primal edges and runs a graph traversal. It checks that the partition is
nontrivial and that each deleted edge has exactly one endpoint on each side.
It then closes each side by identifying its two cut-edge labels, normalizes
the PD, and revalidates spherical one-component structure. A trace stores the
local PD and the chosen edges and crossing partition. Applying the procedure
recursively yields visible diagram factors, not necessarily prime knots.

For a valid spherical projection, a dual two-cycle gives the cut directly;
the primal checks are conservative safeguards. Building the dual and checking
one cut is linear graph work. Repeated splitting has at most $n-1$ cuts and
polynomial overhead, including revalidation and trace construction. A
conservative $O(n^3)$ operation bound covers the Python implementation without
assuming optimal disjoint-set or balanced splitting behavior. A diagram of a
composite knot can hide every such splitting curve; this procedure does not
solve general prime decomposition.

\subsection{Why unreduced ranks can be halved degree by degree}
The characteristic-two splitting is standard in Khovanov theory
\cite{Shumakovitch,Wigderson}. Here is an algebraic proof in the conventions
needed by the program. On a resolution with $r$ circles, write its vector space
as $A^{\otimes r}$, $A=\F[x]/(x^2)$. Fix a basepoint on the knot, let $X$ be
multiplication by $x$ on its circle, and let
\[
 \nu=\sum_{j=1}^r\partial_{x_j},\qquad
 \partial_x(1)=0,\quad\partial_x(x)=1.
\]
Both commute with the differential. For $X$ this is the marked Frobenius-module
relation. For $\nu$ it follows by checking
\[
 \partial m=m(\partial\otimes1+1\otimes\partial),\qquad
 (\partial\otimes1+1\otimes\partial)\Delta=\Delta\partial
\]
on the basis $1,x$, where
\[
 m(x,x)=0,\quad m(1,x)=m(x,1)=x,\quad m(1,1)=1,
 \quad\Delta(1)=1\otimes x+x\otimes1,\quad\Delta(x)=x\otimes x.
\]
Also $X^2=0$ and $\nu X+X\nu=1$.

\begin{proposition}\label{prop:split}
After forgetting quantum grading, the unreduced complex splits into two
isomorphic copies of the reduced complex in the same cube degrees.
Consequently every unreduced homological rank is even.
\end{proposition}
\begin{proof}
Take the reduced complex to be $\im X$. The identity
$\nu X+X\nu=1$ decomposes every chain element into a member of $\im X$ and
a member of $\nu\im X$. If $y=\nu z\in\im X$ with $z\in\im X$, then
$0=Xy=X\nu z=z$, hence $y=0$. Thus
$C=\im X\oplus\nu\im X$. The maps $\nu$ and $X$ restrict to mutually
inverse chain maps between these summands. Their quantum shifts are irrelevant
here; neither changes cube degree. Taking homology proves the claim.
\end{proof}

\subsection{Tensor products and homological convolution}
For a marked knot diagram define the reduced rank polynomial
\[
 R_D(z)=\sum_h\dim_{\F}\rKh^h(D;\F)z^h,
\]
using the package's unnormalized cube degree $h$.

\begin{proposition}\label{prop:product}
For visibly joined diagrams,
\[
 R_{D_1\#D_2}(z)=R_{D_1}(z)R_{D_2}(z).
\]
In particular their total reduced ranks multiply.
\end{proposition}
\begin{proof}
A resolution of the joined diagram is a pair of resolutions of the factors.
Their marked circles become a single marked circle. Label that joined circle
by $x$ in the reduced complex, and retain all labels on the other circles.
This gives a basis bijection with the tensor product of the two reduced
complexes. It respects the local merge and split maps: a crossing change lies
in one factor, and its marked and unmarked label rules are unchanged by the
joining. Thus the differential is $d_1\otimes1+1\otimes d_2$ over $\F$,
with cube degrees adding. Over a field a finite complex splits as its homology
with zero differential plus contractible two-term summands. Tensoring shows
that only the product of the homology summands contributes to homology, which
is exactly the stated coefficient convolution.
\end{proof}
The program scans each factor, divides its per-degree unreduced ranks by two,
convolves the reduced polynomials, and doubles the final coefficients. It never
constructs one Python object for every generator of the joined homology. For
recognition it can stop at the first nontrivial factor; all factors must be
trivial to return $\unk$. The cap on live objects applies separately to the
scans, not to the potentially enormous integer rank returned at the end.

\subsection{An actual family-level asymptotic improvement}
\begin{theorem}\label{thm:family}
Fix a bounded-size knot diagram $D$ with reduced rank $r>1$. On explicitly
joined diagrams $D^{\#k}$ that expose their two-edge splitting curves, the new
factorwise rank algorithm runs in polynomial time in $k$. An algorithm which
materializes the final unreduced minimal complex must use at least $2r^k$
objects. There are such families of bounded scan width.
\end{theorem}
\begin{proof}
The splitting step and its validations are polynomial in $n=k|D|$. Each factor
scan has constant size and cost. Each reduced factor polynomial has fixed
degree; after $j$ factors the product has $O(j)$ coefficients, each with
$O(j)$ bits. Sequential multiplication by a fixed-degree factor is polynomial
in $k$ even with elementary integer arithmetic. Proposition~\ref{prop:product}
gives total reduced rank $r^k$, and Proposition~\ref{prop:split} gives
unreduced rank $2r^k$. A closed minimal complex with zero differential has one
object for each of these dimensions. Finally, scanning the fixed factors
successively keeps only a constant-size factor boundary and the constant
number of joining strands, giving bounded width independently of $k$.
\end{proof}
This is a representation improvement for specified families, not a general
polynomial unknot algorithm. It also refutes a width-only bound for the
\emph{multiplicities} of the explicit Khovanov complex. A small number of
boundary matching types does not bound how many copies of each type survive.

For a particularly relevant family, use the supplied 11-crossing Conway
diagram $C$. Independent dense computation gives reduced rank 33 and exact
Alexander polynomial one. The chosen splice creates no new decreasing R1/R2
move; the supplied diagram itself has none. Its iterated sums are nontrivial,
so no descending traversal exists. Their Alexander polynomial is still one.
The original recognition pipeline therefore reaches its explicit homology
backend, whose final size is $2\cdot33^k$. The new factorwise recognizer handles
this family in polynomial time. This conclusion does not rely on a lucky
Jones specialization of the whole sum; factoring precedes the filters and a
fixed nontrivial factor is enough to decide the sum.

\section{Completeness and general worst-case complexity}\label{sec:complete}
\begin{theorem}\label{thm:correct}
On a validated classical knot diagram, with no user resource limits, the
accelerated recognizer terminates and returns $\unk$ precisely for the unknot.
With limits, every completed verdict is sound; $\unknown$ is not a verdict.
\end{theorem}
\begin{proof}
The original local reductions preserve knot type and a descending diagram is
an unknot. Splitting curves preserve connected-sum structure. The Alexander
and Jones rejection steps are sound by
Propositions~\ref{prop:alexfilter} and~\ref{prop:jonesfilter}. The scan builds
the same finite local complex as the baseline, and its changes preserve the
represented morphisms and homology by Lemmas~\ref{lem:inverse}
and~\ref{lem:cancel}. Factoring preserves ranks by
Proposition~\ref{prop:product}.

Kronheimer--Mrowka prove that reduced rational Khovanov rank one detects the
unknot~\cite{KM}. If the reduced $\F$ rank is one, the universal coefficient
theorem bounds the reduced rational rank above by one. That rational homology
is nonzero: its graded Euler characteristic is the normalized Jones polynomial,
whose value at 1 for a knot is 1. Hence the rational rank is exactly one and
the knot is trivial. Conversely the unknot's reduced $\F$ complex has rank one.
This justifies using $\F$ rank in the recognizer. Products of positive reduced
ranks equal one exactly when every factor rank is one, which also justifies the
factorwise recognition rule.

Every scan has finitely many crossings and objects; every cancellation deletes
two objects. The optional filter cap skips a finite stage rather than aborting
the algorithm. Thus the unlimited process terminates. Limited executions only
stop early with $\unknown$, without changing any algebraic test.
\end{proof}

The implementation admits a conservative $2^{O(n)}$ upper bound on normalized
inputs, not a quasi-polynomial one. Attaching a crossing chooses two smoothings
and creates at most two new closed circles, so prior objects generate at most
eight times as many objects before cancellation. After $n$ steps there are at
most $8^n$ such objects. A morphism has at most $2^{O(n)}$ dot monomials; a
composition, even with naive expansion, has $2^{O(n)}$ cost. There are at most
exponentially many cancellations and polynomially many pairs of incident
entries per cancellation in the exponentially bounded object set. Their
product is still $2^{O(n)}$. The number of generated heap entries obeys the
same type of bound. Coefficient and label bit costs only contribute polynomial
factors to these exponential estimates.

The Alexander stages use polynomially bounded degree and coefficient bit
length, and factorization only splits the crossing set. Thus neither worsens
this bound. The proof is deliberately loose; it is not a prediction of typical
running time. Caches, Markowitz pivots, and greedy order selection have no proved
general quasi-polynomial performance guarantee. Nor do the tests establish a
new complexity lower bound for all possible homology algorithms: the lower
bound in Theorem~\ref{thm:family} concerns an explicit-output representation.

\section{Executed experiments}\label{sec:experiments}
\subsection{Measurement protocol}
The saved measurements use CPython 3.13.5, Linux x86-64, and an Intel Xeon
Platinum 8370C at the reported 2.80 GHz nominal frequency. The environment exposed
five logical CPUs. Each trial starts a fresh subprocess with
\texttt{PYTHONHASHSEED=0}, imports exactly one implementation, validates the
same PD, and then starts the kernel timer. Trials run serially. Startup, imports,
and input validation are excluded from successful kernel timings; external
subprocess caps include them. The machine is a shared execution environment,
so small variations and sub-millisecond differences should not be overread.

\texttt{results/benchmarks.json} contains all 21 cases, exact generated input
PDs, per-trial measurements, returned answers, caps, and environment data.
Ordinary examples and the 256-crossing ordering use three-trial medians. The
36-crossing raw scan, large factored examples other than $C^{\#2}$,
1,024-crossing ordering, and descending microbenchmark use one trial.
A first timed-out baseline trial is recorded as censored; it is not repeatedly
rerun to manufacture a median. The table generator reads this JSON directly.

\subsection{Complete recognition}
\begin{center}
\small
\begin{tabular}{lrrrr}
\toprule
Workload & $n$ & Baseline (s) & New (s) & Ratio \\
\midrule
Trefoil & 3 & 0.000188 & 0.000203 & 0.93 \\
Hard unknot (8) & 8 & 0.003776 & 0.003626 & 1.04 \\
Conway & 11 & 0.092830 & 0.002823 & 32.89 \\
Kinoshita--Terasaka & 11 & 0.082476 & 0.003428 & 24.06 \\
$T(3,5)$ & 10 & 0.000971 & 0.000591 & 1.64 \\
Unknot braid (40) & 40 & 0.014090 & 0.011755 & 1.20 \\
Recorded workload (36) & 36 & 0.014231 & 0.003462 & 4.11 \\
\bottomrule
\end{tabular}
\end{center}

The ratio is baseline time divided by new time; values below one are
regressions. The Conway and Kinoshita--Terasaka cases reach the original homology
backend but are rejected by the new exact Jones evaluation. Their modular
Alexander checks correctly remain inconclusive. The hard eight-crossing unknot
still needs the complete fallback. The trefoil is slightly slower because the
new early-stage overhead is material at a fraction of a millisecond.

\textbf{The 36-crossing recorded timeout is not a difficult complete-recognition
instance.} Both complete pipelines reject it cheaply. The large improvement in
the following subsection concerns an explicitly requested raw homology
computation without the recognition filters or reductions.

\subsection{Unfactored homology backend}
\begin{center}
\small
\begin{tabular}{lrrrr}
\toprule
Diagram & $n$ & Baseline (s) & New (s) & Reduced rank \\
\midrule
Hard unknot (8) & 8 & 0.005712 & 0.004923 & 1 \\
Conway & 11 & 0.086681 & 0.027837 & 33 \\
Kinoshita--Terasaka & 11 & 0.099573 & 0.033056 & 33 \\
$T(3,5)$ & 10 & 0.022543 & 0.021728 & 7 \\
Unknot braid (40) & 40 & 0.019398 & 0.007964 & 1 \\
Recorded workload (36) & 36 & timeout & 2.810384 & 2949 \\
\bottomrule
\end{tabular}
\end{center}

The timed-out entry here is a 12-second external wall cap. A separate,
earlier baseline run allowed 120 seconds of kernel time and stopped at
120.143533 seconds without a result; its exact record and reproduction script
are included. That observation is not a 120-second paired median and is not
used to assign a numeric ratio to the censored table cell. The same workload
completed in the accelerated raw scan in 2.810384 seconds. Its peak was 17,694
objects before elimination and 5,898 after elimination, with maximum boundary
8 and 16,884 cancellations. The final unreduced rank is 5,898, hence reduced
rank 2,949.

An additional scan with $d^2$ checks completed in 3.822706 seconds with the same
rank. An attempted original-backend computation using input crossing order
instead of the default order exhausted its 12-second kernel budget. Thus the
large rank is supported by internal differential checks and the extensive
small-input comparisons below, but this work does \emph{not} claim that the
original backend independently completed that 36-crossing calculation.

The results also show why a single speedup multiplier is inappropriate.
The $T(3,5)$ kernel barely changes, whereas the large cancellation workload
changes dramatically. Cold-cache small examples, not only warm repeated calls,
are included. The original and new ``compositions'' counters have different
bookkeeping conventions, so those raw counters should not be interpreted as
paired operation-count ratios.

\subsection{Factored ranks and preprocessing}
\begin{center}
\small
\begin{tabular}{lrlrrr}
\toprule
Diagram & $n$ & Baseline & New (s) & Reduced rank & Peak objects \\
\midrule
$C^{\# 2}$ & 22 & 5 s cap & 0.035691 & $33^{2}$ & 246 \\
$C^{\# 3}$ & 33 & 5 s cap & 0.077945 & $33^{3}$ & 246 \\
$C^{\# 6}$ & 66 & 5 s cap & 0.096359 & $33^{6}$ & 246 \\
$C^{\# 12}$ & 132 & 5 s cap & 0.176065 & $33^{12}$ & 246 \\
$3_1^{\# 10}$ & 30 & 5 s cap & 0.011318 & $3^{10}$ & 18 \\
\bottomrule
\end{tabular}
\end{center}

The two-factor Conway new time is a three-trial median; the other rows are
single trials. All corresponding baseline first trials hit a 5-second external
wall cap. The listed peak is the maximum number of live objects before
elimination in \emph{any individual factor scan}, not a process memory measure
and not the returned rank. For twelve Conway factors,
\[
 33^{12}=1{,}667{,}889{,}514{,}952{,}984{,}961.
\]
The program returns the full unnormalized homological rank polynomial as well
as this total. The final integer has modest bit length; the original explicit
homology representation would have astronomical dimension. The practical
speedup is consequently explained by an exact change of representation, not
by discarding homology.

\begin{center}
\small
\begin{tabular}{lrrrr}
\toprule
Workload & $n$ & Baseline (s) & New (s) & Ratio \\
\midrule
Greedy order, 256 & 256 & 0.483933 & 0.010180 & 47.54 \\
Greedy order, 1,024 & 1024 & 8.117443 & 0.088788 & 91.43 \\
Descending, 501 & 501 & 0.001404 & 0.001405 & 1.00 \\
\bottomrule
\end{tabular}
\end{center}

The ordering gains agree with the asymptotic change from quadratic to
near-linear work per start. The essentially unchanged descending timing is
reported rather than suppressed. As stressed earlier, a better worst-case
bound for a subroutine need not make every tested input faster.

\subsection{Validation results and what they do not prove}
The final suite consists of 30 passing test groups. Its main independent and
regression checks are as follows.
\begin{itemize}[itemsep=3pt]
\item 120 random one-component braid diagrams with at most eight crossings are
compared to an independent dense reduced cube calculation, both with factoring
and with the unfactored new scanner. Fifty are also compared degree by degree
with the preserved original backend; 24 run intermediate $d^2$ checks.
\item 120 diagrams are evaluated by an independent complete bracket state sum
in each of three rings: the large prime field, $\mathbb F_{101}$, and
$\mathbb Z/15\mathbb Z$. These 360 checks do not use the scanning glue routine.
Additional tests cover mirrors, PD half-turns, crossing permutations,
Reidemeister reductions, braid relations, and Markov stabilizations.
\item 500 random diagrams compare descending-start existence and the greedy
order with the original implementation. Returned descending darts are replayed.
All 128 units on three arcs are checked, as are 200 general morphism
compositions, cache-result mutation, modular Alexander minors, invalid budgets,
Jones cap fall-through, and connected-sum degrees and traces.
\end{itemize}
The unmodified original 19 test groups also pass in a separate process.
Independent dense computations on the 11-crossing Conway and Kinoshita--Terasaka
diagrams give reduced rank 33 in both cases; these larger reference runs are
recorded separately. The tests are deterministic, with seeds in their source.
They provide substantial evidence against implementation mistakes but do not
constitute a formal proof of the Python program. In particular, no Lean
formalization or independently checkable succinct homology certificate has
been produced.

\section{Reproduction, interfaces, and limitations}
From the distribution root, no installation is needed:
\begin{lstlisting}
python -m fastunknot recognize examples/conway.json
python -m fastunknot recognize examples/hard_unknot_8.json
python -m fastunknot khovanov examples/baseline_timeout36.json --no-factor
python -m fastunknot khovanov examples/conway_sum_12.json
python -m unittest discover -s tests -v
python -m unittest discover -s baseline/tests -v
python tools/benchmark.py --hard-timeout 12
\end{lstlisting}
The optional installation command is \texttt{python -m pip install .}. Runtime
dependencies are limited to the standard library. The code targets Python 3.10
or later, but the delivered experiments tested only the stated 3.13.5
environment. The PDF requires an ordinary LaTeX installation; its generated
tables and source are included.

\texttt{recognize} returns $\unk$, $\knot$, or $\unknown$ and records the
selected method. To disable both Alexander stages use
\texttt{use\_alexander=False}; Jones is controlled separately by
\texttt{use\_jones=False}. The CLI counterparts are \texttt{--no-alexander}
and \texttt{--no-jones}. A standalone Jones call explicitly labels equality
\texttt{INCONCLUSIVE}. Khovanov \texttt{order=} is validated as a permutation
and disables factoring, so an explicit order is never silently replaced by
factor-local orders. The raw scan can also be selected with
\texttt{factor=False} or \texttt{--no-factor}.

Time checks are cooperative. They now cover more loops, including the expensive
optional $d^2$ diagnostics, but cannot interrupt a single long Python arithmetic
or geometry operation. A hard service deadline must use an external process
limit. The object cap is not a global memory cap. Likewise, entry-bounded caches
do not impose a fixed byte ceiling. Numerical evidence and replay traces are
helpful audit information, not a substitute for a separately implemented
certificate verifier.

The primary remaining performance limitation is irreducible intermediate
homology multiplicity on diagrams without useful visible splits. Better
pivot heuristics, compiled kernels, or additional one-sided invariants could
improve constants, but do not by themselves supply the missing logarithmic
hierarchy-depth and compressed-state bounds of Section~\ref{sec:quasi}.
A general quasi-polynomial successor should implement and verify those
three-dimensional contracts, rather than relabel the present scanner.

\appendix
\section{Auditable workload and source map}
The archived raw-backend timeout example is the closure of the following braid
on five strands:
\begin{lstlisting}
[-4,-3,4,4,-3,2,1,1,-1,-1,3,-4,-1,3,2,2,-2,-1,
 -3,2,3,-3,-4,-4,-3,3,1,3,-2,-4,-4,3,-3,3,-4,-2]
\end{lstlisting}
The full final unnormalized unreduced degree counts are
\[
\begin{array}{c|rrrrrrrrr}
h&10&11&12&13&14&15&16&17&18\\\hline
\dim&8&40&112&230&386&552&698&784&792\\[2pt]
h&19&20&21&22&23&24&25&26&27\\\hline
\dim&722&588&428&278&158&78&32&10&2
\end{array}
\]
and sum to 5,898. Degree shifts are not normalized by crossing signs; the total
rank is unaffected by this choice.

The central implementation locations are
\texttt{scan.py} (local algebra, sparse cancellation, order selection, factored
rank wrapper), \texttt{filters.py} (modular determinant and bracket DP),
\texttt{decompose.py} (dual cuts and factor traces), \texttt{simplify.py}
(linear descending test), \texttt{diagram.py} (separate oriented and relation
signs), and \texttt{recognize.py} (sound portfolio control flow).
The original code remains available for differential inspection. The patch
changes only the original package's source files; the standalone archive also
provides the additional validation and benchmark artifacts.

\section{Rebuilding the report and checking results}
\begin{lstlisting}
python tools/report_tables.py
cd paper
pdflatex -interaction=nonstopmode -halt-on-error report.tex
pdflatex -interaction=nonstopmode -halt-on-error report.tex
\end{lstlisting}
\texttt{results/unit-tests.txt} and
\texttt{results/baseline-unit-tests.txt} contain the executed suite outputs;
\texttt{results/independent-large-checks.json} contains the larger independent
reference runs and the checked 36-crossing scan. The supplementary 120-second
baseline observation is in \texttt{results/baseline-hard-120s.json}.
An interrupted benchmark can be resumed with
\texttt{--resume --max-cases 2}; completed measurements are retained.
Successful timings are not comparable to external wall caps without accounting
for the different timing scopes. No checksum files, third-party font files,
or copies of the cited research papers are needed to use this distribution.

\clearpage
\begin{thebibliography}{9}
\raggedright
\bibitem{L2026}
M. Lackenby, \emph{Incompressible surfaces, hierarchies and unknot recognition},
arXiv:2607.23350v1, 25 July 2026, especially Sections 9--10.
\url{https://arxiv.org/abs/2607.23350}.
\bibitem{Ltalk}
M. Lackenby, \emph{Unknot recognition in quasi-polynomial time},
talk slides, 109 physical PDF pages; supplied in the input archive.
\url{https://people.maths.ox.ac.uk/lackenby/quasipolynomial-talk.pdf}.
\bibitem{BN}
D. Bar--Natan, \emph{Fast Khovanov homology computations},
J. Knot Theory Ramifications \textbf{16} (2007), 243--255;
arXiv:math/0606318. \url{https://arxiv.org/abs/math/0606318}.
\bibitem{KM}
P. B. Kronheimer and T. S. Mrowka, \emph{Khovanov homology is an
unknot-detector}, Publ. Math. Inst. Hautes \'Etudes Sci. \textbf{113} (2011),
97--208; arXiv:1005.4346. \url{https://arxiv.org/abs/1005.4346}.
\bibitem{Shumakovitch}
A. Shumakovitch, \emph{Torsion of the Khovanov homology},
arXiv:math/0405474. \url{https://arxiv.org/abs/math/0405474}.
\bibitem{Wigderson}
Y. Wigderson, \emph{The Bar--Natan theory splits}, arXiv:1508.05982.
\url{https://arxiv.org/abs/1508.05982}.
\bibitem{Kauffman}
L. H. Kauffman, \emph{State models and the Jones polynomial},
Topology \textbf{26} (1987), 395--407.
\url{https://doi.org/10.1016/0040-9383(87)90009-7}.
\end{thebibliography}
\end{document}
